How can we picture the connections among a large group of objects, such as information, transportation, or social networks?  In this chapter, we study Graph Theory, which is used throughout the mathematical sciences, natural sciences, and social sciences to model such relationships. An especially powerful aspect of graph models is that we can draw pictures representing graphs, as we did in the Handshakes Puzzle (\Cref{act_handshakes}).  

We define a (combinatorial) graph, explore models, introduce key vocabulary, and discuss representations of graphs in \Cref{sec_graph_models} {\secgraphmodels} and \Cref{sec_std_graphs} {\secstdgraphs}. 
Next, in \Cref{sec_color_graphs} {\seccolorgraphs} we introduce graph coloring and applications that can be recast as coloring problems.  Then, in \Cref{sec_subgraph_isomorphic} {\secsubgraphiso} and \Cref{sec_classify_graphs} {\secclassifygraphs} we discuss how to prove that two drawings of graphs represent the same graph (or not) and how to classify all graphs meeting a given set of criteria.
